\chapter{Matrix Multiplication Analysis}

\section{Motivation}

Large Language Model (LLM) inference workloads are dominated by matrix multiplication operations that form the core of attention and feed-forward layers. However, the overall behavior of such workloads on GPUs is influenced by an interplay of compute, memory, and communication characteristics. Before analyzing end-to-end inference pipelines, it is important to develop a clear understanding of how fundamental GPU metrics behave in a controlled and interpretable setting.

To this end, we begin with the tiled matrix multiplication workload—a simplified kernel that captures the essential compute and memory access patterns seen in LLM inference. The goal here is not to optimize matrix multiplication itself, but rather to use it as a microbenchmark to study key GPU metrics, identify relevant performance indicators, and familiarize ourselves with profiling tools that will later be applied to real inference runs.

We employ two primary NVIDIA tools for this study: \textbf{Nsight Systems (NSYS)} and the \textbf{Data Center GPU Management Interface (DCGMI)}. NSYS provides fine-grained timeline traces, kernel-level statistics, and PCIe throughput data, while DCGMI offers continuous runtime counters such as SM activity, tensor utilization, and memory bandwidth usage. Together, these tools allow us to correlate different phases of GPU execution—data transfers, kernel computation, and synchronization—with corresponding hardware metrics.

By first examining matrix multiplication in isolation, we build intuition about how metrics such as SM occupancy, memory bandwidth utilization, and tensor activity reflect GPU efficiency. This understanding serves as a foundation for interpreting the more complex, multi-kernel behavior of LLM inference workloads analyzed later.

\section{Algorithm}
The CUDA kernel \texttt{matmul\_tiled} computes the product of two tiles and accumulates the result in the output tile. Each thread block computes a $B_s \times B_s$ submatrix of the output tile.

\begin{algorithm}[H]
\caption{Tiled Matrix Multiplication using CUDA with Host-Level Tiling}
\label{alg:cuda_matmul}
\begin{algorithmic}[1]
\Require Matrices $A, B \in \mathbb{R}^{N \times N}$, block size $B_s$, tile size $T_s$
\Ensure $C = A \times B$

\State Allocate host matrices $h_A, h_B, h_C$

\Comment{Process in tiles at host level}
\For{$i = 0$ to $N-1$ \textbf{step} $T_s$}
    \For{$j = 0$ to $N-1$ \textbf{step} $T_s$}
        \State Allocate device submatrix $d_C$ of size $T_s \times T_s$
        \State Initialize $d_C \gets 0$
        \For{$k = 0$ to $N-1$ \textbf{step} $T_s$}
            \State Allocate $d_A, d_B$ on device ($T_s \times T_s$)
            \State Copy $T_s \times T_s$ submatrices:
            \State \hspace{1em} $d_A \gets h_A[i : i+T_s,\, k : k+T_s]$
            \State \hspace{1em} $d_B \gets h_B[k : k+T_s,\, j : j+T_s]$
            
            \State Launch CUDA kernel:
            \State \hspace{1em} $\text{matmul\_tiled}<<< \lceil T_s/B_s \rceil,\, \lceil T_s/B_s \rceil >>>(d_A, d_B, d_C, T_s)$
            
            \State Synchronize device
            \State Free $d_A, d_B$
        \EndFor
        \State Copy $d_C$ back to host:
        \State \hspace{1em} $h_C[i : i+T_s,\, j : j+T_s] \gets d_C$
        \State Free $d_C$
    \EndFor
\EndFor
\State \Return $h_C$
\end{algorithmic}
\end{algorithm}



\section{Nsight Systems GPU Metrics}
\begin{table}[h]
\centering
\begin{tabular}{|c|p{11cm}|}
\hline
\textbf{Metric} & \textbf{Description} \\ \hline
GPC Clock Frequency & GPU core clock (Graphics Processing Cluster). Determines SM execution speed. \\ \hline
SYS Clock Frequency & System clock domain of the GPU, affecting memory controllers, PCIe, and copy engines. \\ \hline
GR Active & Fraction of time the graphics engine is active (rendering or compute dispatch). \\ \hline
Sync Compute In Flight & Number of synchronous compute kernels currently executing (blocking). \\ \hline
Async Compute In Flight & Number of asynchronous compute kernels currently executing in parallel (non-blocking). \\ \hline
SM Active & Fraction of time at least one warp is active on any Streaming Multiprocessor (SM). \\ \hline
SM Issue Instructions & Fraction of cycles where SMs issued at least one instruction. Indicates pipeline utilization. \\ \hline
Tensor Active Instructions & Fraction of executed instructions that used Tensor Cores (e.g., matrix multiply). \\ \hline
SM Warp Occupancy & Ratio of active warps to the maximum possible per SM. Measures parallel utilization. \\ \hline
DRAM Bandwidth & Achieved global memory bandwidth (GB/s). Indicates memory throughput and potential bottlenecks. \\ \hline
PCIe RX/TX Throughput & Host-to-Device (RX) and Device-to-Host (TX) data transfer rates over PCIe. \\ \hline
PCIe Read/Write BAR1 Requests & CPU requests to GPU memory aperture (BAR1). Reads = CPU fetching from GPU, Writes = CPU writing to GPU. \\ \hline
\end{tabular}
\caption{Nsight Systems GPU Metrics}
\label{tab:nsys_metrics}
\end{table}


Below is an example of NSYS GPU metrics captured during a matrix multiplication run.

\begin{figure}[H]
    \centering
    \includegraphics[width=0.9\textwidth]{\FIGDIR/nsys_matmul.png}
    \caption{NSYS GPU Metrics for a matrix multiplication run}
\end{figure}


\begin{itemize}
    \item As shown below by the NVTX ranges, the matrix multiplication involves various phases: host to device copy, kernel computation, and device to host copy.
    \item The compute metrics - compute in flight, SM Active, SM Occupancy reach high values during the kernel execution phase, indicating good GPU utilization.
    \item The PCIE bandwidth metrics spike during the copy phases, showing data transfer activity. 
\end{itemize}


\section{Analysis of GPU Utilization using DCGMI}
DCGMI - NVIDIA's Data Center GPU Management Interface - also provides insights into GPU performance metrics. 

A few important metrics to consider are \cite{nvidiadcgmi}:
\begin{table}[h!]
\centering
\begin{tabular}{|c||p{12cm}|}
\hline
\textbf{Metric} & \textbf{Description} \\
\hline
\hline
Graphic Engine Activity & The fraction of time any portion of the graphics or compute engines were active. \\
\hline
SM Activity & The fraction of time at least one warp was active on a multiprocessor, averaged over all multiprocessors. Note that “active” does not necessarily mean a warp is actively computing. For instance, warps waiting on memory requests are considered active. \\
\hline
SM Occupancy & The fraction of resident warps on a multiprocessor, relative to the maximum number of concurrent warps supported on a multiprocessor. The value represents an average over a time interval and is not an instantaneous value. Higher occupancy does not necessarily indicate better GPU usage. \\
\hline
Tensor Activity & The fraction of cycles the tensor (HMMA / IMMA) pipe was active. The value represents an average over a time interval and is not an instantaneous value. Higher values indicate higher utilization of the Tensor Cores. \\
\hline
Memory BW Utilization & The fraction of cycles where data was sent to or received from device memory. The value represents an average over a time interval and is not an instantaneous value. \\
\hline
NVLink Bandwidth & The rate of data transmitted / received over NVLink, not including protocol headers, in bytes per second. \\
\hline
PCIe Bandwidth & The rate of data transmitted / received over the PCIe bus, including both protocol headers and data payloads, in bytes per second. \\
\hline
\hline
\end{tabular}
\caption{DCGMI GPU Metrics}
\label{tab:dcgmi_metrics}
\end{table}


In the figures below, we analyze the GPU utilization metrics (primarily focussed on compute) across different matrix sizes for our tiled matrix multiplication implementation (tile size set equal to the matrix size).


\begin{figure}[H]
    \centering
    \includegraphics[width=\textwidth]{\RESULTDIR/matmul_compute.png}
    \caption{SM Active and SM Occupancy across different matrix sizes}
\end{figure}
